\setcounter{section}{18}

  \zwischenueberschrift{Soal latihan}

\inputaufgabe
{}
{

Misalkan $L,M,N$ adalah
\definitionsverweis {manifold-manifold Riemann}{}{.}
Tunjukkan sifat-sifat berikut.
\aufzaehlungdrei{Identitas
\maabbdisp {\operatorname{Id}_{  M  }} { M } { M
} {}
adalah suatu
\definitionsverweis {isometri}{}{.}
}{Jika
\maabb {\varphi} { L } { M
} {}
adalah suatu isometri, maka
\definitionsverweis {pemetaan invers}{}{}
\maabb {\varphi^{-1}} { M } { L
} {}
juga merupakan suatu isometri.
}{Jika
\maabb {\varphi} { L } { M
} {}
dan
\maabb {\psi} { M } { N
} {}
adalah isometri, maka $\psi \circ \varphi$ juga merupakan suatu isometri.
}

}
{} {}

\inputaufgabe
{}
{

Misalkan $M$ adalah suatu
\definitionsverweis {manifold Riemann}{}{.}
Tunjukkan bahwa himpunan semua
\definitionsverweis {isometri}{}{}
pada $M$ membentuk suatu
\definitionsverweis {grup}{}{.}

}
{} {}

\inputaufgabe
{}
{

Misalkan
\maabbeledisp {\varphi} {\R} {S^1 \subseteq \R^2
} { t } { \begin{pmatrix}  \cos  t  \\ \sin  t    \end{pmatrix}
} {.}
Tunjukkan bahwa $\varphi$ adalah suatu
\definitionsverweis {isometri lokal}{}{}
antara
\definitionsverweis {manifold-manifold Riemann}{}{}
jika $\R$ dan $\R^2$ dilengkapi dengan
\definitionsverweis {struktur Euklides}{}{}
dan $S^1$ dilengkapi dengan struktur Riemann terinduksi.

}
{} {}

\inputaufgabe
{}
{

Misalkan
\maabb {\varphi} {\R^n } { \R^n
} {}
adalah suatu
\definitionsverweis {pemetaan linear}{}{,}
dan misalkan $\R^n$ dilengkapi dengan
\definitionsverweis {struktur Euklides}{}{}
beserta
\definitionsverweis {struktur Riemann}{}{}
yang diinduksikannya. Tunjukkan bahwa $\varphi$ merupakan suatu
\definitionsverweis {isometri linear}{}{}
jika dan hanya jika $\varphi$ merupakan suatu
\definitionsverweis {isometri}{}{}
terhadap struktur Riemann tersebut.

}
{} {}

\inputaufgabe
{}
{

Misalkan

\mavergleichskette
{\vergleichskette
{ I
}
{ \subseteq }{ \R
}
{  }{
}
{    }{
}
{   }{
}
}
{}{}{}
adalah suatu
\definitionsverweis {interval terbuka}{}{}
dan
\maabbdisp {\gamma} { I } { \R^n
} {}
adalah suatu
\definitionsverweis {kurva yang terdiferensialkan secara kontinu}{}{}
yang injektif. Kita memandang $I$ maupun $\R^n$ sebagai
\definitionsverweis {manifold Riemann}{}{}
melalui
\definitionsverweis {struktur Euklides}{}{}
masing-masing. Andaikan bahwa

\mavergleichskette
{\vergleichskette
{ \gamma(I)
}
{ \subseteq }{ U
}
{ \subseteq }{ \R^n
}
{    }{
}
{   }{
}
}
{}{}{}
dengan $U$
\definitionsverweis {terbuka}{}{}
merupakan suatu
\definitionsverweis {submanifold tertutup}{}{.}
Tunjukkan bahwa pemetaan

\mavergleichskettedisp
{\vergleichskette
{ \gamma
}
{ :} { I
}
{ \longrightarrow} { \gamma(I) \subseteq \R^n
}
{ } {
}
{ } {
}
}
{}{}{}
merupakan suatu
\definitionsverweis {isometri}{}{}
jika dan hanya jika $\gamma$
\definitionsverweis {berparametrisasi panjang busur}{}{.}

}
{} {}

\inputaufgabe
{}
{

Misalkan $L_1,L_2,M_1,M_2$ adalah
\definitionsverweis {manifold-manifold Riemann}{}{}
dan misalkan
\maabb {\varphi_1} {L_1 } { M_1
} {}
serta
\maabb {\varphi_2} {L_2 } { M_2
} {}
adalah
\definitionsverweis {isometri}{}{.}
Tunjukkan bahwa
\maabbdisp {\varphi_1 \times \varphi_2} {L_1 \times L_2 } { M_1 \times M_2
} {}
juga merupakan suatu isometri.

}
{} {}

\inputaufgabe
{}
{

Kita meninjau
\definitionsverweis {difeomorfisme}{}{}
\maabbeledisp {\varphi} {S^1 \times \R_+} { \Complex \setminus \{0\}  =  \R^2 \setminus \{(0,0)\}
} {(s,t) } { st
} {,}
dengan silinder $S^1 \times \R_+$ dilengkapi
\definitionsverweis {struktur produk Riemann}{}{.}
Uraikan struktur Riemann pada
\mathl{\R^2 \setminus \{(0,0)\}}{}
yang membuat $\varphi$ menjadi suatu
\definitionsverweis {isometri}{}{.}

}
{} {}

\inputaufgabegibtloesung
{}
{

Misalkan
\mathkor {} {r} {dan} {R} {}
adalah bilangan real positif dengan

\mavergleichskette
{\vergleichskette
{ 0
}
{ < }{ r
}
{ < }{ R
}
{    }{
}
{   }{
}
}
{}{}{,}
dan kita meninjau torus
\zusatzklammer {tertanam} {} {}

\mavergleichskettedisp
{\vergleichskette
{ T
}
{ =} { \left\{ (x,y,z) \in \R^3  \mid   \left(  \sqrt{x^2+y^2} -R  \right)^2 +z^2  = r^2         \right\}
}
{ \subseteq} { \R^3
}
{ } {
}
{ } {
}
}
{}{}{}
dengan
\definitionsverweis {struktur Riemann terinduksi}{}{.}
Tunjukkan bahwa untuk setiap

\mavergleichskette
{\vergleichskette
{ \alpha
}
{ \in }{ [0,2 \pi [
}
{  }{
}
{    }{
}
{   }{
}
}
{}{}{}
pemetaan
\maabbeledisp {\Psi_\alpha} {\R^3} {\R^3
} { \left(  x   , \,   y  , \,  z                 \right) } { \left(  x \cos \alpha - y \sin  \alpha   , \,   x \sin \alpha + y \cos  \alpha  , \,  z                 \right)
} {,}
menginduksi suatu
\definitionsverweis {isometri}{}{}
dari $T$ ke dirinya sendiri.

}
{} {}

\inputaufgabe
{}
{

Hitung panjang lintasan linear dari
\mathl{\begin{pmatrix}  0  \\ 1   \end{pmatrix}}{} ke
\mathl{\begin{pmatrix}  1  \\ 1   \end{pmatrix}}{} dalam
\definitionsverweis {setengah bidang Poincaré}{}{.}

}
{} {}

\inputaufgabe
{}
{

Hitung luas persegi panjang dalam
\definitionsverweis {setengah bidang Poincaré}{}{}
yang diberikan oleh titik-titik sudut
\mathl{(1,1),\, (5,1),\, (1,3),\, (5,3)}{}{.}

}
{} {}

\inputaufgabegibtloesung
{}
{

Tunjukkan bahwa
\definitionsverweis {pemetaan-pemetaan}{}{}
\maabbeledisp {\varphi} { {\mathbb H} } { U \left(  0,1  \right)
} { z} { { \frac{  z-  { \mathrm i}  }{ z + { \mathrm i}  } }
} {,}
dan
\maabbeledisp {\psi} {  U \left(  0,1  \right)  } { {\mathbb H}
} {w} { { \frac{   (w+1)  { \mathrm i}  }{ 1-w  } }
} {,}
memberikan suatu
\definitionsverweis {pemetaan biholomorfik}{}{}
antara
\definitionsverweis {setengah bidang atas}{}{}
${\mathbb H}$ dan cakram satuan terbuka $U \left(  0,1  \right)$.

}
{} {}

\inputaufgabe
{}
{

Tentukan
\definitionsverweis {turunan}{}{}
dari
\definitionsverweis {pemetaan-pemetaan}{}{}
\maabbeledisp {\varphi} { {\mathbb H} } { U \left(  0,1  \right)
} { z} { { \frac{  z-  { \mathrm i}  }{ z + { \mathrm i}  } }
} {,}
dan
\maabbeledisp {\psi} {  U \left(  0,1  \right)  } { {\mathbb H}
} {w} { { \frac{   (w+1)  { \mathrm i}  }{ 1-w  } }
} {.}

}
{} {}

\inputaufgabegibtloesung
{}
{

Uraikan
\definitionsverweis {pemetaan}{}{}
\maabbeledisp {\psi} { \Complex \setminus \{1\} } { \Complex
} {w} { { \frac{   (w+1)  { \mathrm i}  }{ 1-w  } }
} {}
dalam koordinat real.

}
{} {}

\inputaufgabegibtloesung
{}
{

Tentukan
\definitionsverweis {turunan parsial}{}{}
dari
\definitionsverweis {pemetaan}{}{}
\maabbeledisp {\psi} { \R^2 \setminus \{  (1,0 )\} } { \R^2
} {  \begin{pmatrix} a \\b   \end{pmatrix} } {  { \frac{  1 }{ -2a+1+a^2+b^2 } }  \begin{pmatrix} -2b \\1-a^2-b^2   \end{pmatrix}
} {.}

}
{} {}

\inputaufgabe
{}
{

Tentukan titik-titik potong garis yang melalui
\mathl{\begin{pmatrix}  0  \\ 0\\  -1   \end{pmatrix}}{} dan suatu titik
\mathl{\begin{pmatrix}  a  \\ b \\  0   \end{pmatrix}}{} dengan

\mavergleichskette
{\vergleichskette
{ a^2+b^2
}
{ < }{ 1
}
{  }{
}
{    }{
}
{   }{
}
}
{}{}{}
dengan hiperboloid dua lembar; lihat
Contoh 18.11.

}
{} {}

Misalkan
\mathkor {} {L} {dan} {M} {}
adalah dua
$C^k$-\definitionsverweis {manifold}{}{.}
Suatu
\definitionsverweis {pemetaan kontinu}{}{}
\maabbdisp {\varphi} { L } { M
} {}
disebut suatu
\definitionswortpraemath {C^k}{ difeomorfisme lokal }{,}
jika untuk setiap titik

\mavergleichskette
{\vergleichskette
{ P
}
{ \in }{ L
}
{  }{
}
{    }{
}
{   }{
}
}
{}{}{}
terdapat suatu
\definitionsverweis {lingkungan terbuka}{}{}

\mavergleichskette
{\vergleichskette
{ P
}
{ \in }{ U
}
{ \subseteq }{ L
}
{    }{
}
{   }{
}
}
{}{}{}
sedemikian sehingga $\varphi  \vert_U$ menginduksi suatu
\definitionsverweis {difeomorfisme}{}{}
antara $U$ dan $\varphi(U)$.

  \zwischenueberschrift{Soal untuk dikumpulkan}

\inputaufgabe
{}
{

Kita meninjau permukaan elipsoid

\mavergleichskettedisp
{\vergleichskette
{ E
}
{ =} { \left\{ (x,y,z)  \mid   2x^2+3y^2+5z^2 = 1        \right\}
}
{ \subset} { \R^3
}
{ } {
}
{ } {
}
}
{}{}{}
dengan
\definitionsverweis {struktur Riemann}{}{}
yang diinduksi oleh ruang sekitarnya. Tentukan
\definitionsverweis {isometri linear}{}{}
pada $\R^3$ yang menginduksi suatu
\definitionsverweis {isometri}{}{}
pada $E$.

}
{} {}

\inputaufgabe
{4}
{

Misalkan $L,M$ adalah
\definitionsverweis {manifold terdiferensialkan}{}{,}
misalkan
\maabbdisp {\varphi} { L } { M
} {}
adalah suatu
\definitionsverweis {difeomorfisme lokal}{}{,}
dan misalkan $M$ adalah suatu
\definitionsverweis {manifold Riemann}{}{.}
Tunjukkan bahwa terdapat tepat satu struktur Riemann pada $L$ yang membuat $\varphi$ menjadi suatu
\definitionsverweis {isometri lokal}{}{.}

}
{} {}

\inputaufgabe
{3}
{

Berikan suatu contoh
\definitionsverweis {manifold Riemann}{}{}
$M$ dan suatu
\definitionsverweis {difeomorfisme}{}{}
\maabbdisp {\varphi} { M } { M
} {,}
yang bukan suatu
\definitionsverweis {isometri}{}{.}

}
{} {}

\inputaufgabe
{4}
{

Hitung panjang busur pada setengah lingkaran atas berjari-jari $\sqrt{2}$ dari
\mathl{\begin{pmatrix}  1  \\ 1   \end{pmatrix}}{} ke
\mathl{\begin{pmatrix}  -1 \\ 1   \end{pmatrix}}{} dalam
\definitionsverweis {setengah bidang Poincaré}{}{.}

}
{} {}

\inputaufgabe
{4}
{

Hitung luas cakram berpusat di
\mathl{\begin{pmatrix}  2  \\ 2   \end{pmatrix}}{} dan berjari-jari $1$ dalam
\definitionsverweis {setengah bidang Poincaré}{}{.}

}
{} {}

\inputaufgabe
{4}
{

Uraikan suatu
\definitionsverweis {isometri}{}{}
eksplisit antara
\definitionsverweis {setengah bidang Poincaré}{}{}
dan lembar atas hiperboloid dua lembar.

}
{} {}
